%=========== ΛΥΣΗ - ΘΕΜΑ Β_92 ===========
\begin{thema}{Β}
Στο παρακάτω αναβαθμισμένο σχήμα έχει προστεθεί διδακτικός σχολιασμός. \textbf{ΠΡΟΣΟΧΗ}: Υπενθυμίζεται πως δεν πρόκειται για το γράφημα της $f$, αλλά για το \textbf{γράφημα της παραγώγου $f'$}. Οπότε μας ενδιαφέρει κυρίως το **πρόσημό της** και δευτερευόντως η δική της μονοτονία!
\begin{center}
\begin{tikzpicture}
    \begin{axis}[
        axis lines = middle,
        axis line style = {->, >=latex},
        grid = major,
        grid style = {very thin, gray!30},
        xlabel = {$x$},
        ylabel = {$y$},
        xtick = {-2,-1,1,2,3,4,5},
        ytick = {-4,-3,-2,-1,1,2,3},
        every axis x label/.style = {at={(current axis.right of origin)}, anchor=west},
        every axis y label/.style = {at={(current axis.above origin)}, anchor=south},
        width = 10cm, height = 8cm, samples = 100,
        xmin=-2.5, xmax=4.5, ymin=-4.5, ymax=3,
        xtick distance=1, ytick distance=1
    ]
        % Η Παράγωγος f'
        \addplot[very thick, \xrwma, domain=-1.5:3.5] {x^2 - 2*x - 3};
        
        % Σημεία οπτικοποίησης προσήμου της f'
        \node[anchor=south, color=blue, font=\large\bfseries] at (axis cs:-1.2, 0.5) {$f'(x)>0$};
        \node[anchor=south, color=blue, font=\large\bfseries] at (axis cs:3.2, 0.5) {$f'(x)>0$};
        \node[anchor=north, color=red, font=\large\bfseries] at (axis cs:1, -1) {$f'(x)<0$};
        
        % Ρίζες της f' (Ακρότατα της f)
        \addplot[mark=*, mark options={fill=black, scale=1.2}] coordinates {(-1,0)};
        \node[anchor=south east, font=\footnotesize] at (axis cs:-1, 0.2) {Tοπ. Mέγιστο της $f$};
        
        \addplot[mark=*, mark options={fill=black, scale=1.2}] coordinates {(3,0)};
        \node[anchor=south west, font=\footnotesize] at (axis cs:3, 0.2) {Tοπ. Ελάχιστο της $f$};
        
        % Ακρότατο της f' (Σημείο Καμπής της f)
        \addplot[mark=*, mark options={fill=red, scale=1.2}] coordinates {(1,-4)};
        \node[anchor=south west, font=\footnotesize] at (axis cs:1.2, -4) {Αλλαγή κυρτότητας $f$ ($f''=0$)};
        
        % Σημείο f'(0) για το Δ ερώτημα
        \addplot[mark=*, mark options={fill=purple, scale=1.2}] coordinates {(0,-3)};
        \node[anchor=north east, font=\footnotesize, color=purple] at (axis cs:0, -3) {$\lambda = f'(0) = -3$};
        
        \node[below left] at (axis cs:0,0) {$0$};
    \end{axis}
\end{tikzpicture}
\end{center}

Με βάση το εμπλουτισμένο σχήμα:
\begin{erwthma}

\item Για τη μελέτη της μονοτονίας της $f$, ελέγχουμε αποκλειστικά το **πρόσημο** της $f'$: πού βρίσκεται πάνω ή κάτω από τον άξονα $x'x$.
\begin{itemize}
    \item $f'(x) > 0$: Η καμπύλη είναι "πάνω" από τον άξονα στα διαστήματα $(-\infty, -1)$ και $(3, +\infty)$.
    \item $f'(x) < 0$: Η καμπύλη είναι "κάτω" από τον άξονα στο διάστημα $(-1, 3)$.
\end{itemize}
Συνεπώς, υπολογίζοντας και τη συνέχεια:
Η $f$ είναι \textbf{γνησίως αύξουσα} στα διαστήματα \textbf{$(-\infty, -1]$} και \textbf{$[3, +\infty)$}.
Η $f$ είναι \textbf{γνησίως φθίνουσα} στο διάστημα \textbf{$[-1, 3]$}.

\item Τα τοπικά ακρότατα της $f$ συναντώνται στις θέσεις όπου η παράγωγός της μηδενίζεται και αλλάζει πρόσημο.
\begin{itemize}
    \item Στο $x_1 = -1$, η παρούσα συνάρτηση $f'$ μηδενίζεται $f'(-1)=0$ και αλλάζει πρόσημο από $(+)$ σε $(-)$. Συνεπώς, στο \textbf{$x=-1$} η $f$ παρουσιάζει \textbf{τοπικό μέγιστο}.
    \item Στο $x_2 = 3$, η παρούσα συνάρτηση $f'$ μηδενίζεται $f'(3)=0$ και αλλάζει πρόσημο από $(-)$ σε $(+)$. Συνεπώς, στο \textbf{$x=3$} η $f$ παρουσιάζει \textbf{τοπικό ελάχιστο}.
\end{itemize}

\item Για να μελετήσουμε την κυρτότητα της $f$, στρεφόμαστε στη \textbf{μονοτονία της παραγώγου $f'$}, διότι η κυρτότητα εξαρτάται από τη δεύτερη παράγωγο $f''$ (που γεωμετρικά αντιστοιχεί στην $1η$ παράγωγο και στο ρυθμό μεταβολής του ίδιου αυτού σχήματος).
\begin{itemize}
    \item Η συνάρτηση $f'$ είναι γνησίως φθίνουσα στο $(-\infty, 1]$. Αυτό σημαίνει ότι $f''(x) \le 0$, και επομένως η  $f$ είναι \textbf{κοίλη} (στρέφει τα κοίλα κάτω) στο $(-\infty, 1]$.
    \item Η συνάρτηση $f'$ είναι γνησίως αύξουσα στο $[1, +\infty)$. Αυτό σημαίνει ότι $f''(x) \ge 0$, και επομένως η $f$ είναι \textbf{κυρτή} (στρέφει τα κοίλα άνω) στο $[1, +\infty)$.
\end{itemize}
Εφόσον η $f$ αλλάζει κυρτότητα στο σημείο $x_3=1$ και διαθέτει εφαπτομένη (ως παραγωγίσιμη παντού), το σημείο **$Α(1, f(1))$** αποτελεί μοναδικό \textbf{σημείο καμπής} της $C_f$.

\item Αναζητούμε την εφαπτομένη της γραφικής παράστασης $C_f$ (και όχι της $C_{f'}$) στο σημείο με $x_0 = 0$. Η εξίσωση της εφαπτομένης δίνεται από τον τύπο:
\[ y - f(0) = f'(0)(x - 0) \]
Γνωρίζουμε ότι το σημείο επαφής έχει τεταγμένη $f(0) = 2$.
Όσο για την κλίση της εφαπτομένης, αυτή ταυτίζεται με την τιμή $f'(0)$. Διαβάζοντας την τεταγμένη επί της δοσμένης γραφικής παράστασης (της $f'$) για τετμημένη $x=0$, εντοπίζουμε ότι \textbf{$f'(0) = -3$}.
Προβαίνοντας σε αντικατάσταση των δεδομένων, διαμορφώνουμε:
\[ y - 2 = -3(x - 0) \Leftrightarrow y = -3x + 2 \]
Άρα εν τέλει, $\varepsilon: y = -3x + 2$.

\end{erwthma}
\end{thema}
%=========== ΤΕΛΟΣ ΛΥΣΗΣ - ΘΕΜΑ Β_92 ===========
